% BHU PYQ -- Manually transcribed & error-corrected
\documentclass[11pt,a4paper]{article}

\usepackage[a4paper,top=2.5cm,bottom=2.5cm,left=2.8cm,right=2.8cm,headheight=14pt]{geometry}
\usepackage{amsmath,amssymb}
\usepackage[T1]{fontenc}
\usepackage[utf8]{inputenc}
\usepackage{lmodern}
\usepackage{microtype}
\usepackage{fancyhdr}
\usepackage{enumitem}
\usepackage{hyperref}
\usepackage{lastpage}
\usepackage{parskip}

\hypersetup{colorlinks=true,urlcolor=black,linkcolor=black,
  pdftitle={CHB-502: Inorganic Chemistry-III -- BHU PYQ},pdfauthor={Banaras Hindu University}}

\pagestyle{fancy}\fancyhf{}
\renewcommand{\headrulewidth}{0.4pt}\renewcommand{\footrulewidth}{0.4pt}
\fancyhead[L]{\small   \textit{Banaras Hindu University}}
\fancyhead[R]{\small   \textit{CHB-502: Inorganic Chemistry-III}}
\fancyfoot[C]{\small Page  \thepage\ of \pageref{LastPage}}


\newlist{parts}{enumerate}{2}
\setlist[parts,1]{label=(\alph*),leftmargin=*,itemsep=5pt,topsep=3pt}
\setlist[parts,2]{label=(\roman*),leftmargin=*,itemsep=3pt,topsep=2pt}

\newcommand{\pts}[1]{\hfill   \textit{[#1]}}


\begin{document}


\begin{center}
  {\large\bfseries BANARAS HINDU UNIVERSITY}\\[2pt]
  {
\normalsize B.Sc. (Hons.) Semester V Examination, 2016-17}\\[6pt]
  \rule{0.8\linewidth}{0.5pt}\\[6pt]
  {\large\bfseries Paper No. CHB-502: Inorganic Chemistry-III}\\[4pt]
  {
\normalsize Time: 3 Hours\hfill Maximum Marks: 70}
\end{center}
\rule{\linewidth}{0.4pt}
\small



\noindent   \textit{Instructions: Answer five questions in total, including Question~1 which is compulsory.}

\normalsize
\bigskip




\noindent   \textbf{Question 1.} Answer all parts of the following: \pts{5 $ \times$ 2 = 10}

\begin{parts}
  \item Why is $[ \text{Co(NH}_3)_6]^{3+}$ diamagnetic, whereas $[ \text{CoF}_6]^{3-}$ is paramagnetic?
  \item Differentiate between thermodynamic stability and kinetic stability of complexes.
  \item What is the trans-effect? Give one example.
  \item Draw the molecular structure of ferrocene.
  \item Write the IUPAC name and structure of Ziegler-Natta catalyst.
\end{parts}

\medskip\rule{\linewidth}{0.2pt}




\noindent   \textbf{Question 2.}

\begin{parts}
  \item Discuss the crystal field splitting of d-orbitals in octahedral and tetrahedral geometries. Derive the relation between their splitting parameters. \pts{8}
  \item Discuss the variation of lattice energies of the bivalent transition metal halides. Explain why the curve shows two maxima (humps) rather than a straight line. \pts{6}
\end{parts}

\medskip\rule{\linewidth}{0.2pt}




\noindent   \textbf{Question 3.}

\begin{parts}
  \item What is the Orgel energy level diagram? Describe the importance of this diagram for distinguishing octahedral and tetrahedral complexes of $ \text{d}^1$ and $ \text{d}^9$ systems. \pts{8}
  \item Define ligand field stabilization energy (LFSE). Calculate the LFSE for $ \text{d}^4$, $ \text{d}^6$, and $ \text{d}^7$ octahedral complexes in both high-spin and low-spin states. \pts{6}
\end{parts}

\medskip\rule{\linewidth}{0.2pt}




\noindent   \textbf{Question 4.}

\begin{parts}
  \item State the mechanisms of substitution reactions in square-planar complexes. How do the charge on the metal center and nature of leaving group affect the rate? \pts{8}
  \item Discuss the homogeneous hydrogenation of alkenes using Wilkinson's catalyst. Write the catalytic cycle. \pts{6}
\end{parts}

\medskip\rule{\linewidth}{0.2pt}




\noindent   \textbf{Question 5.}

\begin{parts}
  \item What is haptacity? Show that the haptacity of a ligand varies from one organometallic compound to another using cyclopentadienyl complexes. \pts{8}
  \item Discuss the Dewar-Chatt-Duncanson model of bonding in metal-alkene complexes, citing Zeise's salt as an example. \pts{6}
\end{parts}

\vfill
\rule{\linewidth}{0.4pt}

\begin{center}\small
  \href{https://bhu-exam.vercel.app/}{Click here to visit our website}\\[4pt]
  \href{https://me-aryan.vercel.app/}{Click here to visit our contributor}\\[4pt]
  \href{https://quashan-me.vercel.app/}{Click here to visit our contributor}
\end{center}

\end{document}
